%% ====================================================================
%%  Chapter 1 — Introduction
%% ====================================================================
\chapter{Introduction}
\label{ch:introduction}

\section{Background and Motivation}
\label{sec:intro:background}

Indoor localization is an essential task for many emerging technologies and applications. Examples of
those applications range from the tracking of assets within large warehouses, the tracking of patients
within a hospital, augmented reality applications, automated robotic systems, and various
location-aware services within large public buildings. Outside of these buildings, the Global
Positioning System (GPS) is typically utilized to determine the location of an object to within a few
meters of its true location. GPS is, however, ineffective within indoor environments; the GPS signals
are reflected and attenuated by the walls and ceilings that surround indoor locations.

Alternative technologies have therefore been used for indoor localization for the past two decades.
Technologies that use WiFi, Bluetooth, ultra-wideband signals, and inertial sensors have all been
explored as alternatives to GPS for indoor localizing of objects or individuals.

WiFi technology has been the focus of much of the research on indoor localization. WiFi signals are of
interest because WiFi infrastructure is already installed throughout most buildings. Furthermore, the
WiFi hardware installed at each access point (AP) provides Channel State Information (CSI). CSI
describes the attenuation and phase of each of the subcarriers that are emitted by each AP. Because the
attenuation of the signals to the WiFi devices is related to the distance between those devices and the
APs, the CSI information can in principle allow an object to be localized to within sub-meter distances
of its actual location.

In order to determine a device's location within a floorplan, DLoc~\cite{ayyalasomayajula2020deep}
utilizes techniques similar to image translation. The CSI from each of the APs is transformed into a
multi-channel image, and a deep learning model with an encoder--decoder structure is trained to read
the CSI from each AP and to output a spatial heatmap indicating the probability of the location of the
object within the building. As a result, DLoc was able to achieve sub-meter accuracy in its
localization results~\cite{ayyalasomayajula2020deep}.

The accuracy that DLoc achieves is, however, at the cost of high model complexity. DLoc contains
16.5~million parameters that are required to perform the localization. Thus, DLoc is not well suited
for deployment upon edge devices such as IoT gateways and smartphones, whose computational and memory
resources are limited.

\section{Problem Definition}
\label{sec:intro:problem}

In order for WiFi localization models to be usefully deployed into indoor environments, the models must
be accurate, robust, and efficient. While DLoc is accurate in localizing objects within a floorplan, it
is not efficient in relation to the resources required for deployment to edge devices. A method of
compressing DLoc such that an accurate, efficient model can be deployed upon edge devices is therefore
required. Such a method of compression comes with several problems that must be resolved.

The high parameter count of models like DLoc implies a high cost for deployment upon edge devices.
Furthermore, it is not known how small a model can be while still achieving accurate localization.

Accuracy is likely to be lost if a model is compressed to reduce its parameter count. An accurate model
that is also compressed for efficiency is therefore desired.

The signals used for WiFi localization are likely to be degraded during their transmission from AP to
WiFi device. Localization models must therefore be robust against signal degradation.

The indoor environments in which the objects to be localized are moving may change over time. The
localization models must therefore be able to generalize to changes in the indoor environment.

The results of training a model are dependent upon the random initialization of its weights. The
localization model must therefore exhibit reproducibility in its results across different random
initializations of its weights.

The problem to be solved is therefore how to replace the DLoc model with a model that is an order of
magnitude smaller yet maintains its localization accuracy, while additionally addressing the robustness
and generalization of the model to changes in the indoor environment.

\section{Project Objectives}
\label{sec:intro:objectives}

The goal of this project is to design, train, and evaluate a family of lightweight models that are
compressed from the DLoc model, and to gain an understanding of the various mechanisms that affect their
accuracy and robustness. The specific objectives of this project are as follows.

\begin{enumerate}[leftmargin=2em]
  \item To compress the DLoc model by training lightweight models (MobileNetV2, a minimal TinyCNN, and a
  state-space Mamba model) to behave similarly to the teacher DLoc model; this allows us to measure the
  accuracy of models that have reduced parameter counts.

  \item To improve the accuracy of the models by using architectural designs (such as U-Net skip
  connections) so that a model with fewer parameters can surpass the accuracy of the teacher DLoc.

  \item To characterize the robustness of the models to signal degradation such as AP dropout, noise,
  signal attenuation, and spatial blur.

  \item To measure the generalization of the models across different sessions of WiFi measurements
  within the indoor environment.

  \item To validate the results of the experiments performed on the models, and to quantify the
  computational efficiency of each of the models tested.
\end{enumerate}

The remainder of this thesis is dedicated to reviewing the relevant literature on indoor localization,
knowledge distillation, efficient deep learning models, and robustness to signal degradation
(Chapter~\ref{ch:literature}). Following that literature review is a description of the DLoc model, the
WILD dataset, and each of the models to be tested in this thesis (Chapter~\ref{ch:methodology}).
Chapter~\ref{ch:implementation} describes the details of how each of the models was implemented.
Chapter~\ref{ch:results} describes the results of testing each of the models, and
Chapter~\ref{ch:conclusion} draws conclusions from those results. Finally,
Chapter~\ref{ch:futurework} proposes future work related to this thesis.
